\section{Proofs and supplementary results for Section~2}
\label{app:structural-proofs}

\begin{proof}[\normalfont\bfseries Proof of Proposition~\ref{thm:cycle}]
    Fix any infinite price path. At time \(t\), the payoff relevant history is the vector
    \[
        s_t=(p_{t-\mathcal A},\ldots,p_t)\in\{p^{L},p^{H}\}^{\mathcal A+1}.
    \]
    Choosing the next price \(p_{t+1}\in\{p^{L},p^{H}\}\) shifts this vector from
    \[
        (p_{t-\mathcal A},\ldots,p_t)
        \quad\text{to}\quad
        (p_{t-\mathcal A+1},\ldots,p_t,p_{t+1}),
    \]
    which is exactly an edge in \(\mathcal G_{\mathcal A}\). Hence every feasible infinite price path induces an infinite walk in \(\mathcal G_{\mathcal A}\). Conversely, every infinite walk in \(\mathcal G_{\mathcal A}\) induces a feasible price path by reading the last coordinate of each visited state.

The node weight \(w(s)\) was defined so that \(w(s_t)\) is exactly the one period revenue at time \(t\). Thus the seller's long run average revenue is the long run average node weight of the corresponding walk. Let \(\lambda^\star\) be the maximum mean weight of a directed cycle in \(\mathcal G_{\mathcal A}\). Any finite prefix of a walk in a finite graph can be decomposed into a residual simple path plus directed cycles. The residual path has length at most the number of nodes, while each extracted cycle has mean weight at most \(\lambda^\star\). Therefore the average weight of a prefix of length \(N\) is at most \(\lambda^\star+O(1/N)\). Taking \(N\to\infty\) shows that no infinite price path can beat \(\lambda^\star\).

    For the reverse inequality, take a directed cycle attaining \(\lambda^\star\) and post its periodic price sequence from time zero. After at most \(\mathcal A\) periods, the fixed initial history has left every memory window and the induced states follow that cycle. The initial discrepancy affects only finitely many bounded rewards, so its contribution to average revenue vanishes. The posted prices are periodic from time zero and their long run average revenue is \(\lambda^\star\). Hence the optimal long run value equals the maximum mean cycle value, and an optimal policy can be chosen periodic.
\end{proof}

\subsection{Revenue reversal symmetry}
\label{sec:reversal}
Reversal preserves revenue because it permutes the promotion counts in windows of each length.
\begin{proposition}[Revenue reversal symmetry]\label{thm:reversal}
Under Assumption~\ref{ass:count-based}, any price cycle \(C=(p_1,\ldots,p_{\mathcal T})\) and its reversal \(C_{\mathrm{rev}}=(p_{\mathcal T},\ldots,p_1)\) have the same average revenue.
\end{proposition}
\begin{proof}
Write \(\widetilde d_t=d_{\mathcal T+1-t}\) for the reversed promotion indicators. \textcolor{black}{For each \(b\),}
\[
 \widetilde N_{b,t}=\sum_{i=0}^b d_{\mathcal T+1-t+i}
 =N_{b,\mathcal T+1-t+b}.
\]
Cyclic summation therefore gives
\(\sum_t F_b(\widetilde N_{b,t})=\sum_t F_b(N_{b,t})\), including when windows exceed the cycle length. Reversal also preserves \(\mathcal T\) and \(m\). Proposition~\ref{thm:exposure-decomposition}, whose proof does not use reversal symmetry, now gives equal cycle revenues. Dividing by \(\mathcal T\) proves the claim.
\end{proof}
Thus a search over cycles can retain one representative from each class related by rotation or reversal.
